% For Joy: please enter the implementation for both MLP modes in the next subsection, change title if required
% (Joy, 4 Nov 2026: done below, in red.)
{\JOY
\subsection{Problem Formulation and Data}
\label{approx_problem}

In this part of the study the price series is not extended into the
future but reproduced: each model receives the whole series of scaled
Close prices \(\{p_1,\ldots,p_T\}\), indexed by the rescaled time
\(t_i=(i-1)/(T-1)\in[0,1]\), and is fitted to it in the mean-squared
sense with a fixed number of parameters. There is no train/test split;
the quantity of interest is the in-sample error as a function of model
size, which isolates the representational capacity of an architecture
from the unpredictability of the target. The comparison is organised
around five parameter budgets, \(N^{*}\in\{15,30,60,120,240\}\); each
model is configured to come as close to the budget as its integer
settings allow, and the actual count \(N\) of every fit is reported.
RMSE and maximum absolute error are recorded in each asset's scaled
\([0,1]\) price units.

Three assets are used: the CBA.AX series of
Section~\ref{data_collection}, and two chosen to differ from it in
character, TSLA (annualised volatility 70\%, maximum drawdown 74\% over
the period) and LLY (volatility 32\%, the strongest sustained trend among
fourteen large-cap candidates screened). Daily data for TSLA and LLY
were obtained from Yahoo Finance for the same window, 2~January 2020 to
31~July 2023 (901 trading days), and prepared exactly as in
Section~\ref{data_collection}, each feature Min--Max scaled to \([0,1]\)
over the full series; only the scaled Close is approximated.

\subsection{KAN Configuration}
\label{approx_kan}

The KANs are built with the PyKAN implementation of the forecasting
study, but several settings of Table~\ref{tab:kan_config} are changed
by the task, and Table~\ref{tab:approx_kan_config} lists every row of
that table with the value used here. The input is the scalar \(t\)
rather than a flattened window, so there is no look-back and no
heuristic hidden-layer size. Two widths are compared. The single-edge
network \texttt{KAN(width=[1,1])} has no hidden layer and its grid size
\(G\in\{6,21,51,111,231\}\) is chosen to meet each budget exactly
(\(N=G+9\)). The hidden-layer network \texttt{KAN(width=[1,\(h\),1])}
keeps the forecasting grid \(G=3\) and sets \(h\in\{1,1,2,5,10\}\)
(\(N\approx24h\); the two smallest budgets cannot be met and are
reported at \(N=24\)). The spline grid is placed on \([0,1]\), the
domain of \(t\), instead of the PyKAN default \([-1,1]\), and the
adaptive grid update that PyKAN applies by default during training is
switched off, so that the parameter count and knot positions are fixed
throughout the fit. Training uses full-batch L-BFGS for a fixed
80 steps with no early stopping, instead of the 10 steps of the
forecasting study, since the fit here is to the whole series with no
risk of overfitting a held-out set; whether this budget is sufficient
for each width is examined together with the results. Each configuration is fitted
from three seeds \(\{0,1,2\}\) rather than one.

\begin{table}[h]
\centering\JOY
\caption{KAN settings in the approximation study, row by row against Table~\ref{tab:kan_config}}
\label{tab:approx_kan_config}
\footnotesize
\begin{tabular}{lll}
\hline
\textbf{Parameter} & \textbf{Forecasting (Table~\ref{tab:kan_config})} & \textbf{Approximation} \\
\hline
Framework & PyTorch + PyKAN & same \\
Architecture & \texttt{[input, \(N_h\), \(H\)]} & \texttt{[1,1]} and \texttt{[1,\(h\),1]} \\
Input representation & flattened window \(L\cdot|\mathcal{F}|\) & scalar \(t\in[0,1]\) \\
Hidden neurons & \(3\times\lfloor|\text{train set}|/50\rfloor\) & none; \(h\in\{1,1,2,5,10\}\) by budget \\
Grid size & 3 & \([1,h,1]\): 3; \([1,1]\): \(G\in\{6,21,51,111,231\}\) \\
Grid range / update & \([-1,1]\) (default), adaptive & \([0,1]\), fixed (no update) \\
\(k\) (B-spline degree) & 3 & same \\
Optimiser & L-BFGS (full-batch) & same \\
Training steps & 10 & 80 \\
Regularisation & PyKAN defaults only & same \\
Random seed & 0 (fixed) & \(\{0,1,2\}\) \\
Stopping criteria & fixed 10 steps, no early stopping & fixed 80 steps, no early stopping \\
\hline
\end{tabular}
\end{table}

\subsection{MLP Configuration}
\label{approx_mlp}

The MLP of Eq.~\eqref{eq:mlp} is implemented in
PyTorch~\cite{paszke2019pytorch} with two hidden layers of width \(w\),
\texttt{tanh} units and a linear output, in two modes: the input is
either the scalar \(t\) (\(N=w^2+4w+1\)) or its Fourier-feature map
\(\gamma(t)\) of Eq.~\eqref{eq:fourier} with \(m\) fixed random
frequencies drawn from \(\mathcal{N}(0,6^2)\) once per seed
(\(N=w^2+(2m+3)w+1\)); the frequencies are not trained and not counted.
Both modes are trained by full-batch Adam~\cite{kingma2015adam}
(learning rate \(10^{-2}\), cosine-annealed) on the mean squared error
for a fixed 20{,}000 steps with no early stopping, retaining the weights
of lowest training loss; PyTorch default initialisation is used, with
no dropout or weight decay, from seeds \(\{0,1,2\}\). Table~\ref{tab:approx_config} lists the settings and
actual parameter counts per budget for all four models.

\begin{table}[h]
\centering\JOY
\caption{Approximation study: settings and actual parameter counts \(N\) per budget}
\label{tab:approx_config}
\small
\begin{tabular}{lcccc}
\hline
\textbf{Budget \(N^{*}\)} & \textbf{KAN [1,1]} & \textbf{KAN [1,\(h\),1]} & \textbf{MLP of \(t\)} & \textbf{Fourier MLP} \\
 & \(G\) (\(N\)) & \(h\) (\(N\)) & \(w\) (\(N\)) & \((m,w)\) (\(N\)) \\
\hline
15  & 6 (15)    & 1 (24)   & 2 (13)   & (2,2) (19)  \\
30  & 21 (30)   & 1 (24)   & 4 (33)   & (2,3) (31)  \\
60  & 51 (60)   & 2 (48)   & 6 (61)   & (4,4) (61)  \\
120 & 111 (120) & 5 (120)  & 9 (118)  & (6,6) (127) \\
240 & 231 (240) & 10 (240) & 14 (253) & (8,9) (253) \\
\hline
\end{tabular}
\end{table}

As in Section~\ref{sec:comparison_protocol}, the optimisers differ in
kind (80 second-order steps against 20{,}000 first-order steps), so
recorded fit times reflect these choices rather than an intrinsic cost
of either architecture; this is taken into account when the results
are discussed.
}
